% ===== 原文 pp.74–76（PDF p.84–86） =====
% [原文 p.74]
周期标量的 Neumann 条件的对偶是规范场的 Dirichlet 条件，反之亦然。下文我们将反复用到这个众所周知的事实。

\subsubsection{DD 条件}

第一种可能是 $t=1$ 时规范扇区与物质扇区同时取 Dirichlet 条件。A-型规范扇区的 Dirichlet 条件映为以 $T^*[2]\C^*$ 为靶空间的 Rozansky-Witten 模型中的一个边界条件，该条件在边界上令 $\sigma=0$，而让复标量 $\tau$ 自由涨落。A-型物质扇区的 Dirichlet 条件则映为 B-型规范理论中的 Neumann 条件。注意 A-型物质扇区的 Dirichlet 条件有两个实参数，分别取值于 $S^1$ 与 $\R$。前一个参数被映为边界 theta 角，即作用量中如下形式的边界项：
\[
  \theta\int_{x_3=0}\frac{F}{2\pi}
  =\theta\int_{x_3=0}\frac{\cF}{2\pi}.
\]
后一个参数是场 $\phi_3$ 的边界值。如第 4 节所讨论的，这两个参数都是不相关的。

% [Tutor Note: p.74 上式右端的花体 $\cF$ 是 B-型规范场的场强，左端 $F$ 是 A-型规范场的场强；文字层把两端都抽成 “F”（见 split/p084.txt），必须按 p.84 图像区分。等式成立是因为对偶把 A-侧的边界 theta 角 $\theta\int F/2\pi$ 映为 B-侧同形式的拓扑项。]

% [Tutor Note: p.74 “as discussed in section 4”（第 4 节）的指代存疑：$A_4$ 与 $\phi_3$ 边界值的不相关性是在 §6.4.2.1（pp.71–72）论证的，而 §6.4 恰是第 6 章的第 4 节。译者按原文直译为“第 4 节”，读者应理解为 §6.4（特别是 §6.4.2.1）。]

如前所述，A-型 3d 规范理论中边界线算子的范畴是 $T^*U(1)$ 的 Fukaya-Floer 范畴。另一方面，Rozansky-Witten 模型中边界线算子的范畴是 $D^b(\mathrm{Coh}(\C^*))$，如 [36] 所解释。这两个范畴由通常的 2d 镜像对称而等价。

让我们回顾 2d 镜像对称在此情形下如何作用于某些对象。$\C^*$ 上的平凡线丛映到底空间 $S^1$ 某个特定点上的纤维。这个特定点使我们能把 $S^1$ 与群流形 $U(1)$ 等同起来。更一般地，我们可以考虑 $\C^*$ 上带有如下形式 $\dbar$-联络的全纯线丛
\[
  \dbar+\ii\lambda\frac{\dd\bar z}{\bar z},\qquad \lambda\in\C .
\]
我们将把这样的线丛记作 $\cL_\lambda$。规范变换可用来消去 $\lambda$ 的虚部，也可以把 $\lambda$ 的实部平移任意整数。因此我们可以把参数 $\lambda$ 看作取值于 $\R/\Z\simeq S^1$。镜像对称把 $\cL_\lambda$ 映到 $T^*U(1)$ 中的一个拉格朗日子流形，即点 $\exp(2\pi\ii\lambda)\in U(1)$ 上的纤维。

% [Tutor Note: p.74 为何 $\lambda$ 的有效取值是 $\R/\Z\simeq S^1$：$\lambda$ 的两个实参数中，$\operatorname{Re}\lambda$ 给出 $\C^*$ 线丛平坦联络的和乐，只在模整数（大规范变换）时有意义；$\operatorname{Im}\lambda$ 只是平凡化的实标度，可被规范变换消去。故 $\lambda$ 的有效空间是 $\R/\Z\simeq S^1$，与镜像侧“$T^*U(1)$ 在 $\exp(2\pi\ii\lambda)$ 处的纤维”逐点对应。]

% [原文 p.75]
把镜像对称作用于 $D^b(\mathrm{Coh}(\C^*))$ 上由导出张量积给出的显然幺半结构，我们就得到 $T^*U(1)$ 的 Fukaya 范畴上的一个幺半结构。$\C^*$ 上的平凡全纯线丛（它在 $D^b(\mathrm{Coh}(\C^*))$ 中充当单位对象）映到 $U(1)$ 的单位元上的拉格朗日纤维。若考虑点 $\exp(2\pi\ii\lambda_1)$、$\exp(2\pi\ii\lambda_2)\in U(1)$ 上的两个拉格朗日纤维，它们的镜像分别是线丛 $\cL_{\lambda_1}$ 与 $\cL_{\lambda_2}$。它们的张量积是线丛 $\cL_{\lambda_1+\lambda_2}$，其镜像是点 $\exp(2\pi\ii(\lambda_1+\lambda_2))\in U(1)$ 上的拉格朗日纤维。显然，这条关于 Fukaya 范畴对象张量积的规则用到了 $U(1)$ 的群结构，即它是卷积型的张量积。

另一类值得考虑的自然拉格朗日子流形是 $T^*U(1)$ 的常值截面，即由方程 $\phi_4=\text{const}$ 给出的子流形。这些子流形是圆，可以带有非平凡的平坦联络。因此这样的 A-膜由 $\R\times U(1)\simeq\C^*$ 的点标号。其镜像对象是 $\C^*$ 上的摩天楼层。支撑在不同点上的两个摩天楼层的导出张量积显然是零对象。一个摩天楼层与自身的导出张量积可以证明同构于该摩天楼层与它平移 $-1$ 之和。也就是说，它是同一点上的摩天楼层，其茎是分次向量空间 $\C[-1]\oplus\C$。应用镜像对称，我们看到 $T^*U(1)$ 的一个截面与自身的张量积必须是同一截面的两份拷贝之和，但其中一份的 Maslov 分次平移了 $-1$。我们不知道如何不借助镜像对称、即通过计算 A-型规范理论中边界线算子的乘积来重现这一结果。

% [Tutor Note: p.75 “摩天楼层与自身的导出张量积是 $\C[-1]\oplus\C$”可用 Koszul 复形核对：在光滑点 $p$ 处 $\cO_p\otimes^{\mathbb{L}}\cO_p\simeq\cO_p\otimes\Lambda^\bullet T_p^*\C^*$；因 $\C^*$ 是复一维的，只有 $\Lambda^0$ 与 $\Lambda^1$ 两个分量，恰对应 $\C[-1]\oplus\C$ 的两项。镜像侧“同一截面的两份拷贝、其中一份 Maslov 分次移 $-1$”与之逐项对应。]

如前所述，A-型物质扇区中边界线算子的范畴是某个相当不寻常的 2d TFT 中膜的范畴，该 TFT 是以 $T^*\R$ 为靶空间的 A-模型的某种修正。前面已经论证过，这个范畴等价于 $D^b_{\C^*}(\mathrm{Coh}(\bullet))$。这与 B-侧一致：在 B-侧，区间上的约化给出一个 B-型 2d 规范理论。

把规范扇区与物质扇区合在一起，我们看到 A-侧的 DD 边界条件映到我们所谓 B-侧的特定（distinguished）边界条件。这样的边界条件对应的边界线算子范畴是 $\C^*$-等变的凝聚层导出范畴 $D^b_{\C^*}(\mathrm{Coh}(\C^*))$，配以其显然的幺半结构。在 A-侧我们得到 $T^*U(1)$ 的 Fukaya-Floer 范畴的一个分级版本：其中拉格朗日子流形上的平坦向量丛带有额外的整数分次，且态射要求关于该分次为次数零。这个分次来自周期标量 $A_4$ 的缠绕数。

我们也可以从 4d 的角度解释这个对偶。事实上，容易看出 A-侧的 DD 边界条件来自 $t=1$ 时的 4d Dirichlet 边界条件，而它在 B-侧的对偶来自 $t=\ii$ 时的 4d Neumann 条件。因此，正如所预期的，电磁对偶交换 4d 中的 Dirichlet 与 Neumann 边界条件。对应于这类 4d 边界条件的面算子可以按如下方式解释：我们切除面算子支撑的一个管状邻域，并在由此产生的边界上施加 4d 边界条件。在 TFT 中，这样的步骤给出一个面算子（即无需取管状邻域厚度趋于零的极限）。

% [原文 p.76]
\subsubsection{NN 条件}

这个条件是 A-侧的特定边界条件，因为规范群在边界上未被破坏，且周期标量 $A_4$ 可以自由地探索整个圆。电磁对偶把它映为 B-型规范理论的 Dirichlet 边界条件，以及以 $\C^*$ 为靶空间的 RW 模型中固定 $\C^*$-值标量 $\tau$、而让 $\sigma$ 自由的边界条件。注意，A-型规范理论中的 Neumann 边界条件与 RW 模型中相应的边界条件都有一个取值于 $\C^*\simeq\R\times U(1)$ 的参数。

让我们比较边界线算子的范畴。A-型规范理论中边界线算子的范畴是凝聚层的上有界导出范畴 $D^b(\mathrm{Coh}(\C[2]))$。RW 模型中边界线算子的范畴也是 $D^b(\mathrm{Coh}(\C[2]))$。A-型物质扇区中边界线算子的范畴是 $T^*S^1$ 的 Fukaya-Floer 范畴。B-型规范扇区中边界线算子的范畴是 $D^b(\mathrm{Coh}(\C^*))$。它们的等价是通常 2d 镜像对称的一个特例。

但还有更多内容：我们预期这些边界线算子的范畴作为幺半范畴也是等价的。对于以 $\C^*$ 为靶空间的 RW 模型与规范群为 $U(1)$ 的 A-型规范理论，这一点可以直接看出。事实上，在两种情形中，边界线算子在范畴中的典型对象都是全纯向量丛的复形，它们可以由边界上某个超联络在拉回向量丛上的 Wilson 线算子来表示。在经典近似下，融合两个这样的边界线算子对应于复形的张量积，而这一结果不可能有量子修正。

比较另一对对偶理论（B-型规范理论与 A-型物质）的幺半结构要更复杂一些。我们不打算在 A-侧做独立计算，而将描述 B-侧的幺半结构，然后说明它在 A-侧对应于什么。

注意，由于 $\C^*$ 是复李群，范畴 $D^b(\mathrm{Coh}(\C^*))$ 有两个自然的幺半结构：导出张量积与卷积型乘积。前者不利用群结构，而后者利用。前者的单位对象是 $\C^*$ 上的全纯函数层，而后者的单位对象是单位点 $1\in\C^*$ 处的摩天楼层。描述 B-侧边界线算子融合的正是后一个幺半结构。事实上，$\C^*$ 上坐标的 3d 含义是联络 $A+\ii\phi$ 沿一个小半圆的和乐，该半圆两端都落在边界上，圆心在边界线算子处（见图 6.1）。摩天楼层对应于这样的边界线算子：该和乐是固定的。特别地，$1\in\C^*$ 处的摩天楼层对应于“不可见的”边界线算子，其和乐为平凡。按定义，这就是边界线算子幺半范畴中的单位对象。

% [Tutor Note: p.76 区分 $D^b(\mathrm{Coh}(\C^*))$ 的两个幺半结构：导出张量积的单位是 $\cO_{\C^*}$（其纤维处处一维，任何点上的摩天楼层与之张量积仍是自身），卷积乘积的单位是 $1\in\C^*$ 处的摩天楼层。书中取后者，因为边界线算子的融合必须与 3d 中沿边界的和乐（即 $\C^*$ 坐标）相容；若误用导出张量积的单位，p.79 与 p.82 的“单位对象不可见”论述会自相矛盾。]

% [Tutor Note: p.76 “不存在量子修正”的依据是：B-型拓扑理论与 A-型规范理论中边界线算子由超联络的和乐给出，其融合只依赖拉回丛上的经典复形结构；拓扑扭转后的理论没有能改变该乘积的量子上同调类。这一点与 §A.1 中“A-型 2d 规范理论的膜范畴为 $\C[2]$ 上凝聚层的有界导出范畴”相容。]
